\section*{第 \ref{ch:g5:air-pressure-weather} 章——大气压与天气}

\begin{solution}{exo:g5:air-pressure-weather:1}
同一个球先充足气称重，再放掉气后称重，
\omterm{def:g3:levers-and-scales:scale}{天平}读数变小——只有\omterm{def:g2:air-around-us:air}{空气}跑了出去，减少的数值就是逃出的\omterm{def:g2:air-around-us:air}{空气}的重量。
这说明\omterm{def:g2:air-around-us:air}{空气}是有\omterm{def:g3:measuring-mass:mass}{质量}的物质。
\end{solution}

\begin{solution}{exo:g5:air-pressure-weather:2}
大气底层受挤压的\omterm{def:g2:air-around-us:air}{空气}对
它所接触到的每一个表面的推力。
这一推力来自我们上方很深的\omterm{def:g2:air-around-us:air}{空气}“海洋”的重量，
并且从各个方向推：向下、向两侧，也向上。
\end{solution}

\begin{solution}{exo:g5:air-pressure-weather:3}
因为各个方向的推力是平衡的：你手掌另一侧的\omterm{def:g2:air-around-us:air}{空气}——
以及你身体内部的\omterm{def:g2:air-around-us:air}{空气}——同样大小地向回推。
只有\emph{不平衡的}推力，比如纸片或被压扁的瓶子所受的推力，
才会显示出它的威力。
\end{solution}

\begin{solution}{exo:g5:air-pressure-weather:4}
水的重量把纸片向下压，大气则把纸片向上托；
大气占了上风，
纸片不掉，水被留在杯中。
\end{solution}

\begin{solution}{exo:g5:air-pressure-weather:5}
先把吸管里的\emph{\omterm{def:g2:air-around-us:air}{空气}}吸走。
大\omterm{def:g5:air-pressure-weather:pressure}{气压}在杯中果汁的液面上，
把果汁沿已排\omterm{def:g2:air-around-us:air}{空气}的吸管向上推：托起果汁的是大气。
\end{solution}

\begin{solution}{exo:g5:air-pressure-weather:6}
把杯子压平，挤走了杯子与瓷砖之间的\omterm{def:g2:air-around-us:air}{空气}。
下面没有\omterm{def:g2:air-around-us:air}{空气}向上推，
大气的推力便把杯子牢牢压在瓷砖上。
\end{solution}

\begin{solution}{exo:g5:air-pressure-weather:7}
大\omterm{def:g5:air-pressure-weather:pressure}{气压}会减小——上方可以压下来的\omterm{def:g2:air-around-us:air}{空气}“海洋”变薄了。
迹象有：快速爬升时耳朵发胀、在稀薄的\omterm{def:g2:air-around-us:air}{空气}中呼吸困难，
以及水开得太早、\omterm{def:g3:temperature-thermometers:temperature}{温度}不够高。
\end{solution}

\begin{solution}{exo:g5:air-pressure-weather:8}
大气的推力，即大\omterm{def:g5:air-pressure-weather:pressure}{气压}。
\omterm{def:g5:air-pressure-weather:pressure}{气压}快速下降预示暴风雨即将来临：应把野餐搬到室内。
\end{solution}

\begin{solution}{exo:g5:air-pressure-weather:9}
在山顶密封的瓶子里装着稀薄的山上\omterm{def:g2:air-around-us:air}{空气}；
到了山谷，外部的压强大于瓶内被困住的压强，
瓶子于是被压扁了。反过来，在山谷密封的瓶子
里装着\emph{较强}的\omterm{def:g2:air-around-us:air}{空气}，随着上升，外部压强减小，
瓶子便被撑得鼓鼓的——
每次开车上山，密封的薯片袋也会如此。
\end{solution}

\begin{solution}{exo:g5:air-pressure-weather:10}
\omterm{def:g2:air-around-us:air}{空气}从\omterm{def:g5:air-pressure-weather:pressure}{气压}高的地方持续流向\omterm{def:g5:air-pressure-weather:pressure}{气压}低的地方，
这种流动就是风。
这样，前面的定义便完整了：风是运动的\omterm{def:g2:air-around-us:air}{空气}，
由\omterm{def:g5:air-pressure-weather:pressure}{气压}差所推动。
\end{solution}

\begin{solution}{exo:g5:air-pressure-weather:11}
随着高度升高，大气对水面的压强减小，水更容易沸腾，\omterm{def:g4:changes-of-state:boiling-point}{沸点}会降到
远低于 \qty{100}{\celsius}。根据沸腾时\omterm{def:g3:temperature-thermometers:temperature}{温度}保持不变的规律，沸水的\omterm{def:g3:temperature-thermometers:temperature}{温度}
不会超过其\omterm{def:g4:changes-of-state:boiling-point}{沸点}：多加的火力只会产生更多水蒸气。因此，
山顶的水沸腾时\omterm{def:g3:temperature-thermometers:temperature}{温度}较低，\emph{并且就停留在这一\omterm{def:g3:temperature-thermometers:temperature}{温度}}：
茶泡得淡，面条煮不烂，
物理规律照样成立。
\end{solution}
